\section*{Chapitre \ref{ch:dv:options-market-making-in-practice} --- Tenue de marché d'options en pratique}

\begin{solution}{exo:dv:options-market-making-in-practice:1}
En l'absence de flux informé, le profit vaut $20\,e^{-w/0.3}w$. Sa dérivée $20\,e^{-w/0.3}(1-w/0.3)$ s'annule en $w=0.3$, où le gain par transaction
et la perte due aux transactions se compensent à la marge.
\end{solution}

\begin{solution}{exo:dv:options-market-making-in-practice:2}
$0.42\times0.115=0.048$ et $0.42\times0.199=0.084$~: 4.8 et 8.4 cents par action, 4.83~\$ et 8.37~\$ par contrat.
\end{solution}

\begin{solution}{exo:dv:options-market-making-in-practice:3}
$2q=65\,894$, $fN=3\,000$, $N+2q=75\,894$~: $65\,894\times3\,000/75\,894=2\,605$.
\end{solution}

\begin{solution}{exo:dv:options-market-making-in-practice:4}
Le gain net vaut $g\,2q\,fN/(N+2q)-cq=aq/(N+2q)-cq$ avec $a=2fNg$. Sa dérivée vaut $aN/(N+2q)^2-c$, nulle en $(N+2q)^2=aN/c$, donc
$q^*=(\sqrt{aN/c}-N)/2$~; la dérivée seconde est négative. Avec $a=288\,000$, $N=10\,000$, $c=0.5$~: $q^*=(75\,895-10\,000)/2=32\,947$.
\end{solution}

\begin{solution}{exo:dv:options-market-making-in-practice:5}
Négocier jusqu'à l'edge revient à passer la plus petite transaction qui rétablit un risque acceptable. Négocier jusqu'au centre suppose une transaction plus importante et ramène la position à un point où le
prochain mouvement aléatoire a autant de chances de nécessiter une transaction que de n'en nécessiter aucune. À l'intérieur de la bande, le risque est accepté par construction, donc payer pour le réduire davantage est un gaspillage.
\end{solution}

\begin{solution}{exo:dv:options-market-making-in-practice:6}
L'erreur vaut le skew multiplié par le mouvement logarithmique depuis l'ajustement, divisé par $\sqrt\tau$. L'écart type de ce mouvement croît comme la racine carrée de l'intervalle, et
l'effet du skew sur la volatilité d'un strike est plus fort pour les échéances courtes, où le smile est plus pentu en \omterm{def:dv:implied-volatility-and-its-surface:surface}{log-moneyness}~: d'où $1/\sqrt\tau$.
\end{solution}

\begin{solution}{exo:dv:options-market-making-in-practice:7}
Les calls assignés aux teneurs de marché suivent une loi hypergéométrique~: $65\,894$ positions vendeuses sur $75\,894$, $72\,894$ exercices. Avec 200\,000 tirages, le profit net moyen
s'élève à 108\,552~\$ avec un écart type de 870~\$ (871~\$ d'après la formule hypergéométrique)~; son premier centile vaut 106\,502~\$. Le risque d'assignation est faible à cette taille.
\end{solution}

\begin{solution}{exo:dv:options-market-making-in-practice:8}
Une seule trajectoire. Sur 20 graines du même modèle jouet, la différence moyenne s'élève à 9\,585~\$ avec un écart type de 11\,597~\$ (une erreur type de 2\,593~\$), et les limites ont fait moins bien
dans 6 cas sur 20. Les limites réduisent le risque d'inventaire et, en moyenne dans ce modèle jouet, améliorent aussi le P\&L, car le flux est unilatéral. Mais un seul mois
n'est qu'un tirage unique.
\end{solution}

\begin{solution}{pb:dv:options-market-making-in-practice:1}
\textbf{1.} Oui~: le dividende (50~\$ par contrat) dépasse la valeur du put (2~\$)~; ne pas exercer coûte 48~\$ par contrat.
\textbf{2.} À taux nuls, l'exercice anticipé d'un call n'apporte rien d'autre que le dividende, que reçoit le détenteur de l'action et que ne reçoit pas le détenteur du call.
\textbf{3.} Plus de la moitié des positions acheteuses qui auraient dû être exercées ne l'ont pas été (1996--2006), ce qui a coûté aux détenteurs plus de 491~millions de dollars.
\textbf{4.} D'autres positions vendeuses, y compris celles des dealers. Leurs positions vendeuses sont assignées moins qu'elles ne le seraient si chaque détenteur exerçait, de sorte qu'elles conservent le dividende
sur la part non assignée.
\textbf{5.} Les teneurs de marché ajoutent $2q$ positions vendeuses qui partagent les assignations avec les vendeurs d'origine. Quand $q$ augmente, leurs positions vendeuses deviennent presque la totalité de
l'encours vendeur à découvert, de sorte que presque toute la non-assignation retombe sur eux.
\textbf{6.} $2q\,fN/(N+2q)$.
\textbf{7.} $fN=3\,000$ contrats, vers lequel on tend quand $q\to\infty$.
\textbf{8.} 0.50~\$~: quatre transactions à 0.10~\$ et deux exercices à 0.05~\$.
\textbf{9.} 32\,947 contrats dans chaque sens~; gain net de 108\,553~\$ (2\,605 non assignés, 125\,026~\$ bruts, 16\,474~\$ de frais).
\textbf{10.} 28\,591~\$ (16\,909 contrats) et 193\,510~\$ (43\,990 contrats).
\textbf{11.} Le gain diminue~; l'écart type issu de l'assignation aléatoire vaut environ 870~\$ à la meilleure taille.
\textbf{12.} Moins d'échecs d'exercice~: le gain baisse, et à la taille prévue les frais peuvent le dépasser (avec 10~\% d'échecs d'exercice, la meilleure taille est deux fois plus petite).
\textbf{13.} Les frais croissent linéairement avec $q$, si bien que la meilleure taille et le profit diminuent~; sans frais, la stratégie capterait l'intégralité des 144\,000~\$ à une taille illimitée.
\textbf{14.} La position en actions issue des calls exercés et tout déséquilibre d'assignation, reportés à l'ouverture de la date de détachement, ainsi que le risque opérationnel sur les instructions d'exercice.
\textbf{15.} De l'historique des exercices de la série lors des dates de détachement passées, de la part des positions ouvertes détenues par des comptes de particuliers, et de son degré de profondeur dans la monnaie.
\textbf{16.} Les détenteurs qui ne parviennent pas à exercer. Les teneurs de marché ne provoquent pas la défaillance, mais ils détournent sa valeur des vendeurs d'origine vers eux-mêmes.
\textbf{17.} Seulement à une taille qui couvre les frais~: avec 10~\% de défaillances, la meilleure transaction rapporte un gain net de 28\,591~\$ après des frais d'environ 8\,500~\$, ce qui reste rentable avec ces chiffres.
\textbf{18.} L'exercice automatique des calls dont le dividende dépasse la valeur temps, ou des frais sur les stratégies de dividende suffisamment élevés pour annuler $q^*$.
\textbf{19.} Avec 30~\% des 10\,000 contrats du public non exercés, un gain de 48~\$ par call non assigné et 0.50~\$ de frais par contrat négocié dans chaque sens~: 108\,553~\$ à la meilleure
taille de 32\,947 contrats (28\,591~\$ à 10~\%, 193\,510~\$ à 50~\%).
\textbf{20.} Des détenteurs d'options qui ne parviennent pas à exercer, collectés par la règle d'assignation au prorata.
\end{solution}

\begin{solution}{iq:dv:options-market-making-in-practice:1}
Exprimée en points de volatilité, puis multipliée par le \omterm{def:dv:greeks-and-the-hedging-pnl:greeks}{vega}~: assez large pour couvrir l'erreur de la \omterm{def:dv:options-market-making-in-practice:theo}{valeur théorique} face aux traders informés, assez étroite pour conserver le flux non informé~; estimée à partir des
mark-outs des transactions passées, par type de contrepartie, et ajustée pour l'inventaire et les limites.
\iqlookfor{erreur de la \omterm{def:dv:options-market-making-in-practice:theo}{valeur théorique}, sélection adverse, \omterm{def:dv:fourier-pricing-and-calibration-engineering:vega}{pondération par le vega}, mark-outs.}
\end{solution}

\begin{solution}{iq:dv:options-market-making-in-practice:2}
Pas à intervalles fixes~: dans une bande autour du delta cible, dont la largeur croît avec le coût de couverture et décroît avec le gamma et l'aversion au risque (demi-largeur proportionnelle
à la racine cubique du coût multipliée par le carré du gamma), en négociant jusqu'à l'edge de la bande. Sur l'exemple du chapitre, une bande procure le même coût qu'une couverture quotidienne, avec
41~\% de bruit en moins sur le P\&L.
\iqlookfor{la bande, et pourquoi elle l'emporte sur un rythme fixe.}
\end{solution}

\begin{solution}{iq:dv:options-market-making-in-practice:3}
À la veille d'une date de détachement, les teneurs de marché négocient de gros volumes compensatoires d'un call profondément dans la monnaie et exercent leurs positions acheteuses. L'assignation au prorata laisse la plupart des
calls non exercés du public assignés à d'autres vendeurs, de sorte que les teneurs de marché conservent le dividende sur leurs positions vendeuses non assignées. Les détenteurs qui ne parviennent pas à exercer
perdent, et les vendeurs d'origine perdent ce qu'ils auraient conservé.
\iqlookfor{assignation au prorata, non-exercice, qui perd.}
\end{solution}

\begin{solution}{iq:dv:options-market-making-in-practice:4}
Rapide~: le forward, ainsi que la \omterm{def:dv:options-market-making-in-practice:theo}{valeur théorique} et le delta de chaque série à partir des paramètres courants de la surface, par des formules fermées et des quantités mises en cache par échéance, avec la règle de
valorisation entre deux recalibrages~; les cotations sont mises à jour dès que la \omterm{def:dv:options-market-making-in-practice:theo}{valeur théorique} varie de plus d'un seuil. Lent~: recalibrages de la surface (à intervalles réguliers et sur événement), contrôles de limites, paramètres de largeur issus des
mark-outs. Les cotations sont retirées pendant qu'un recalibrage s'exécute.
\iqlookfor{séparation du chemin critique et des recalibrages, recalibrages déclenchés par événement.}
\end{solution}

\begin{solution}{iq:dv:options-market-making-in-practice:5}
Les \omterm{def:dv:greeks-and-the-hedging-pnl:greeks}{vegas} compensatoires sur des échéances différentes sont exposés aux variations de la structure par terme, que le \omterm{def:dv:greeks-and-the-hedging-pnl:greeks}{vega} total masque. À proximité d'une limite, le desk décale ses cotations pour attirer
des transactions compensatoires, se couvre avec des options cotées dans le \omterm{def:dv:options-market-making-in-practice:bucket}{bucket de vega}, et ne retire le côté qui franchirait la limite qu'à la limite.
\iqlookfor{risque de structure par terme, décalage des cotations avant retrait.}
\end{solution}

\begin{solution}{iq:dv:options-market-making-in-practice:6}
C'est le risque de pin~: une assignation incertaine. Réduire la position avant la clôture (racheter le strike ou négocier le pin), dimensionner la position en actions du week-end en fonction de
l'assignation attendue, garder une capacité de couverture à l'ouverture du lundi~; connaître les heures limites d'exercice et le seuil d'exercice automatique.
\iqlookfor{incertitude de l'assignation et action avant la clôture.}
\end{solution}
